\chapter{Gaussian Realization at \texorpdfstring{$c_*$}{c*}}
\label{ch:generative-gaussian}

This chapter asks what the directed law of
\Cref{ch:generative} becomes when every kernel is linear Gaussian, and what
the blocks of the resulting precision mean. The central tool is one
Schur-complement lemma (\Cref{lem:gauss-schur}): the marginal precision of a
set of directions is a Schur complement, dominated in the Loewner order by the
precision restricted to those directions, with equality exactly under
\(\Lambda\)-orthogonality. With it the chapter computes the assembled
precision of a directed forest (\Cref{prop:gauss-assembled-precision-blocks}),
proves frame covariance of every kernel under the common frame
(\Cref{prop:gen-lg-kernel-covariance}), and states the Gaussian forms of the
M-coordinate, the comparison lift, and the linear record of selective
communication. The associated-bundle morphism \(\widetilde\Phi\) of
\Cref{ch:geometry} is realized locally by the law map
\eqref{eq:gen-lg-induced-law-map} of a root kernel, of which the bridge matrix
\(\mathsf B_i\) is one parameter. The interaction family of \Cref{ch:gaussian}, the
recognition costs of \Cref{ch:restrictions}, the Fisher charts of
\Cref{ch:infogeometry}, and the coarse maps of
\Cref{ch:gaussian-coarsegraining} all start here.

The Gaussian is a tractable realization, not the ambient theory, and no matrix
identity of this chapter is asserted for a non-Gaussian belief fiber. The
common frame ties the gauge choices of the two channels but does not identify
the dimensions \(K\) and \(d_m\), make \(\mathsf B_i\) invertible, turn \(\Phi\) and
\(\widetilde\Phi\) into inverses, or add a reverse generative factor.
\status{NOT-CLAIMED}

\section{Frames and sectors}
\label{sec:gen-lg-frame-group}

\hypothesisheading{Represented frame group of the Gaussian realization}{hyp:gen-lg-frame-group}
Write \(\Gamma_b:=\rho_b(G)\subseteq\GL(K)\) and
\(\Gamma_m:=\rho_m(G)\subseteq\GL(d_m)\) for the represented frame groups of
the two channels. Where a result names this hypothesis, it adopts the
specialization \(\Gamma_b=\GL^+(K)\), realized for example by
\(G=\GL^+(K)\) with \(\rho_b\) its defining representation. Identities that
use only the frame-change laws, such as the kernel covariance of
\Cref{prop:gen-lg-kernel-covariance} and the congruence covariance of
precisions and pencils, hold for every \(\Gamma_b\). Three families of
statements use the specialization: the gain-absorption and orientation
statements of \Cref{sec:gauss-hypothesis}, which need exactly the
orientation-preserving frames; the frame dependence of ordinary eigenvalues in
\Cref{sec:ig-gauge}, which needs \(\Gamma_b\supseteq\R_{>0}I_K\); and the
rank-only classification of \Cref{prop:grg-residual-fiber-label}, which needs
\(\Gamma_b\supseteq\GL^+(K)\). The hypothesis excludes orientation-reversing
represented frames and compact structure groups, whose represented frames lie
in \(\operatorname{O}(K)\) after a choice of invariant inner product.

\propositionheading{Compact represented frames}{prop:gen-lg-compact-frames}
Let \(\Gamma_b\subseteq\operatorname{O}(K)\). A block reframing
\(\mathbf P=\operatorname{diag}(\mathbf P_1,\dots,\mathbf P_N)\) with \(\mathbf P_i\in\Gamma_b\) acts on a
precision by the orthogonal similarity \(\mathbf P^{-\top}\Lambda\mathbf P^{-1}=\mathbf P\Lambda\mathbf P^\top\),
so ordinary spectra are frame invariant, as are log determinants of
\(\Lambda\) restricted to a subspace and computed in orthonormal bases. For
\(M\succ0\) and \(\mathbf P_0\in\Gamma_b\) the congruence \(\mathbf P_0^\top M\mathbf P_0\) is an
orthogonal similarity, so the eigenvalues of \(M\) are invariants; when
\(\Gamma_b\) is \(\operatorname{SO}(K)\) or \(\operatorname{O}(K)\), the orbit
of \(M\) is exactly its orthogonal-similarity class, labeled by its \(K\)
eigenvalues. \status{ESTABLISHED}

\noindent\emph{Proof.}
For orthogonal \(\mathbf P\), \(\mathbf P^{-\top}=\mathbf P\). A restricted log determinant in an
orthonormal basis \(E\) becomes one in the orthonormal basis \(\mathbf PE\), because
\((\mathbf PE)^\top(\mathbf P\Lambda\mathbf P^\top)(\mathbf PE)=E^\top\Lambda E\). A
spectral decomposition \(M=ODO^\top\) with \(O\in\operatorname{O}(K)\) can be
taken with \(\det O=1\) by reversing one eigenvector, so \(\operatorname{SO}(K)\)
already reaches every orthogonal similarity. \(\square\)

\subsection{The common frame}
\label{sec:gen-gaussian-diagonal}

Throughout this supplement the belief and model channels of agent \(i\) use
one local frame \(u_i\), the gauge choice of the main manuscript
(\Cref{sec:geo-diagonal}); one passive re-expression \(g_i\in G\) then acts on
samples by
\begin{equation}
k_i\longmapsto\rho_b(g_i)k_i,
\qquad
m_i\longmapsto\rho_m(g_i)m_i.
\label{eq:gen-gaussian-diagonal-action}
\end{equation}
The two channels may still carry inequivalent representations and
independent link families.

The belief-sector precision analyzed in \Cref{ch:gaussian,ch:restrictions}
is that of the belief conditional \(p(k\given m,o,X)\), the fast belief at
frozen model latents in the sense of \Cref{sec:obs-fast-slow}: bridge terms,
offsets, and the model part of each observation mean shift only the
information vector at fixed \(m\), as
\Cref{prop:gauss-global-interaction-condition} computes. The model sector
\(p(m\given k,o,X)\) has the same structure with \(K\) replaced by \(d_m\).

\section{Information form}
\label{sec:gauss-information-form}

Let \(\nu\) be Lebesgue measure on \(\R^n\). The Gaussian family in
information form has the information vector \(\mathsf h\in\R^n\) and the
precision \(\Lambda\in\Sym_{++}^n\), with density
\begin{equation}
q(y\given\mathsf h,\Lambda)
=\exp\left(\mathsf h^\top y-\tfrac12 y^\top\Lambda y-\mathsf A(\mathsf h,\Lambda)\right),
\qquad y\in\R^n.
\label{eq:gauss-density}
\end{equation}
The sans-serif \(\mathsf A\) is the log normalizer; the self terms of
\Cref{ch:gaussian} are the \(K\times K\) matrices \(A_i\).

\subsection{Coordinate charts}
\label{sec:ig-expfam}

\definitionheading{Gaussian coordinate charts}{def:gauss-coordinate-charts}
Pair two symmetric matrices by the trace form \(\langle M,M'\rangle=\Tr(M^\top M')\) and use the
chart of \(\Sym^n\) whose coordinates are the diagonal entries and the
off-diagonal entries scaled by \(\sqrt2\), in which the Euclidean inner product
of coordinate vectors is this pairing. With the statistic
\(\mathsf T(y)=(y,yy^\top)\), the natural parameter is
\begin{equation}
\lambda=(\lambda_1,\lambda_2)=\left(\mathsf h,-\tfrac12\Lambda\right),
\label{eq:ig-natural}
\end{equation}
the information-form chart is \((\mathsf h,\Lambda)\), the expectation chart
is \(\boldsymbol\mu=(\mu,M_2)\) with \(M_2=\E_q[YY^\top]\), and the moment
chart is \((\mu,\Sigma)\). The moment chart is not Legendre dual to the
natural chart, because \(\Sigma=M_2-\mu\mu^\top\). With the unscaled
half-vectorization every off-diagonal derivative would acquire a factor two;
all matrix blocks in this supplement are written in the scaled chart.

\propositionheading{Log normalizer, moments, and dual charts}{prop:gauss-log-normalizer-moments}
(i) For every \((\mathsf h,\Lambda)\in\R^n\times\Sym_{++}^n\) the density
\eqref{eq:gauss-density} is normalizable, with
\begin{equation}
\mathsf A(\mathsf h,\Lambda)
=\tfrac12\mathsf h^\top\Lambda^{-1}\mathsf h-\tfrac12\log\det\Lambda+\tfrac n2\log2\pi,
\label{eq:gauss-lognorm}
\end{equation}
and
\begin{equation}
d\mathsf A=\mu^\top d\mathsf h-\tfrac12\Tr\left[\left(\Sigma+\mu\mu^\top\right)d\Lambda\right],
\qquad
\mu=\Lambda^{-1}\mathsf h,
\qquad
\Sigma=\Lambda^{-1},
\label{eq:gauss-moments}
\end{equation}
so \(q(\cdot\given\mathsf h,\Lambda)\) has mean \(\mu\) and covariance
\(\Sigma\); the Hessian of \(\mathsf A\) in \(\mathsf h\) at fixed
\(\Lambda\), which is the Fisher information in \(\mathsf h\), is \(\Sigma\),
and the Fisher information in \(\mu\) at fixed \(\Sigma\) is \(\Lambda\).
(ii) In the natural chart the family is the regular
minimal canonical family
\begin{equation}
q(y)=\exp\left(\langle\lambda,\mathsf T(y)\rangle-\mathsf A(\lambda)\right),
\qquad \lambda\in\R^n\times\{\lambda_2\prec0\},
\label{eq:ig-density}
\end{equation}
with
\begin{equation}
\mathsf A(\lambda)
=-\tfrac14\lambda_1^\top\lambda_2^{-1}\lambda_1-\tfrac12\log\det(-2\lambda_2)+\tfrac n2\log2\pi,
\label{eq:ig-lognormalizer}
\end{equation}
strictly convex on its open convex natural domain, and with mean parameter
\begin{equation}
\boldsymbol\mu=\nabla_\lambda\mathsf A=\left(\mu,M_2\right),
\qquad M_2=\Sigma+\mu\mu^\top.
\label{eq:ig-expectation}
\end{equation}
(iii) The Fisher metric in the natural chart is
\begin{equation}
g_\lambda=\nabla^2_\lambda\mathsf A(\lambda)=\operatorname{Cov}_q\left(\mathsf T(Y)\right),
\label{eq:ig-fisher-hessian}
\end{equation}
and in the expectation chart it is
\begin{equation}
g_{\boldsymbol\mu}
=\left(\frac{\partial\lambda}{\partial\boldsymbol\mu}\right)^{\top}g_\lambda
\left(\frac{\partial\lambda}{\partial\boldsymbol\mu}\right)
=g_\lambda^{-1},
\label{eq:ig-dual-inverse}
\end{equation}
the Hessian of \(\mathsf A^\ast\), whose value at \(\boldsymbol\mu\) is the
negative differential entropy of \(q\). \status{ESTABLISHED}

\noindent\emph{Proof.}
(i) Completing the square,
\begin{equation}
\mathsf h^\top y-\tfrac12 y^\top\Lambda y
=-\tfrac12\left(y-\Lambda^{-1}\mathsf h\right)^\top\Lambda\left(y-\Lambda^{-1}\mathsf h\right)
+\tfrac12\mathsf h^\top\Lambda^{-1}\mathsf h,
\label{eq:gauss-complete-square}
\end{equation}
and the Gaussian integral \((2\pi)^{n/2}(\det\Lambda)^{-1/2}\) give
\eqref{eq:gauss-lognorm}; differentiating termwise with
\(d\Lambda^{-1}=-\Lambda^{-1}d\Lambda\Lambda^{-1}\) and
\(d\log\det\Lambda=\Tr(\Lambda^{-1}d\Lambda)\) gives \eqref{eq:gauss-moments}.
At fixed \(\Lambda\) the family is canonical in \(\mathsf h\) with statistic
\(y\), so its Fisher information is the covariance \(\Sigma\); at fixed
\(\Sigma\) the score in \(\mu\) is \(\Lambda(y-\mu)\), with covariance
\(\Lambda\Sigma\Lambda=\Lambda\). (ii) In the scaled chart \(\langle\lambda,\mathsf T(y)\rangle
=\mathsf h^\top y-\tfrac12y^\top\Lambda y\), and substituting
\(\mathsf h=\lambda_1\), \(\Lambda=-2\lambda_2\) into \eqref{eq:gauss-lognorm}
gives \eqref{eq:ig-lognormalizer}. The natural domain is open, and the family
is minimal, because a nonzero affine function of \((y,yy^\top)\) is a nonzero
polynomial of degree at most two and vanishes only on a Lebesgue-null set.
\Cref{prop:exp-cumulants-fisher} then gives strict convexity,
\(\nabla_\lambda\mathsf A=\E\mathsf T\), which is \eqref{eq:ig-expectation},
and \eqref{eq:ig-fisher-hessian}. Since
\(\boldsymbol\mu=\nabla\mathsf A(\lambda)\), the Jacobian
\(\partial\boldsymbol\mu/\partial\lambda\) is \(g_\lambda\), which is positive
definite, so the chain rule gives \eqref{eq:ig-dual-inverse}; by Legendre
duality \(\nabla\mathsf A^\ast(\boldsymbol\mu)=\lambda\), so
\(g_\lambda^{-1}=\nabla^2\mathsf A^\ast\), and
\(\mathsf A^\ast(\boldsymbol\mu)=\langle\lambda,\boldsymbol\mu\rangle-\mathsf A(\lambda)
=\E_q\log q\). \(\square\)

By (i) the precision is at once a coupling matrix and the Fisher metric of
the mean, and the chapters below use both readings.

\section{Conditional reading and the Schur lemma}
\label{sec:gauss-conditional-reading}

Let \(\mathcal P\) be a partition of the \(n\) coordinates into blocks of
dimensions \(n_b\), and write \(\Lambda_{bc}\) for the corresponding blocks.

\propositionheading{Conditional reading of the precision}{prop:gauss-conditional-precision}
Let \(Y\sim\mathcal N(\mu,\Lambda^{-1})\) with \(\Lambda\succ0\). (i) For any
block \(b\),
\begin{equation}
\E\left[Y_b\given Y_{-b}=y_{-b}\right]
=\mu_b-\Lambda_{bb}^{-1}\sum_{c\neq b}\Lambda_{bc}\left(y_c-\mu_c\right),
\qquad
\operatorname{Cov}\left(Y_b\given Y_{-b}\right)=\Lambda_{bb}^{-1}.
\label{eq:gauss-conditional}
\end{equation}
(ii) For distinct blocks \(b\) and \(c\), \(Y_b\) and \(Y_c\) are conditionally
independent given all remaining coordinates if and only if
\(\Lambda_{bc}=0\). (iii) If \(b\) and \(c\) have equal dimension and
\(\Lambda_{bc}=-W_{bc}\) with \(W_{bc}\succeq0\), the displacement
\(\Lambda_{bb}^{-1}W_{bc}\delta_c\) of the conditional mean of \(Y_b\) caused
by \(\delta_c=y_c-\mu_c\) points toward \(\delta_c\) in the \(\Lambda_{bb}\)
inner product, since
\(\delta_c^\top\Lambda_{bb}(\Lambda_{bb}^{-1}W_{bc}\delta_c)=\delta_c^\top W_{bc}\delta_c\geq0\);
in the Euclidean inner product it can point away: for
\(\Lambda_{bb}=\operatorname{diag}(1,10)\),
\(W_{bc}=\begin{psmallmatrix}1&1\\1&1\end{psmallmatrix}\),
\(\Lambda_{cc}=10I_2\), whose joint matrix has leading minors \(1,10,89,780\),
and \(\delta_c=(1,-2)\), the displacement is \((-1,-\tfrac1{10})\) and its
Euclidean pairing with \(\delta_c\) is \(-\tfrac45\). \status{ESTABLISHED}

\noindent\emph{Proof.}
At fixed \(y_{-b}\) the exponent of \eqref{eq:gauss-density} is
\(-\tfrac12y_b^\top\Lambda_{bb}y_b\) plus a linear term, so the conditional law
is Gaussian with precision \(\Lambda_{bb}\), and completing the square around
\(\mu\) gives the conditional mean. Conditioning on the coordinates outside
\(b\cup c\) leaves a Gaussian in \((y_b,y_c)\) with precision
\(\begin{psmallmatrix}\Lambda_{bb}&\Lambda_{bc}\\\Lambda_{cb}&\Lambda_{cc}\end{psmallmatrix}\),
which factorizes exactly when the quadratic form has no cross term. The
numbers in (iii) are direct evaluation. This is the Gaussian case of the
pairwise Markov property \citep[Ch.~5]{Lauritzen1996}; see also
\citet[Ch.~2]{RueHeld2005}. \(\square\)

\lemmaheading{Schur complement}{lem:gauss-schur}
Let \(\Lambda\in\Sym_{++}^n\) and let \([E\ F]\) be an orthogonal matrix with
\(E\in\R^{n\times k}\). (i) The marginal precision of \(E^\top Y\) under
\(\mathcal N(\mu,\Lambda^{-1})\) is the Schur complement
\begin{equation}
\left(E^\top\Lambda^{-1}E\right)^{-1}
=E^\top\Lambda E-E^\top\Lambda F\left(F^\top\Lambda F\right)^{-1}F^\top\Lambda E;
\end{equation}
for the coordinate block \(b\) of a partition this reads
\begin{equation}
\left(\Lambda^{-1}\right)_{bb}
=\left(\Lambda_{bb}-\Lambda_{b,-b}\Lambda_{-b,-b}^{-1}\Lambda_{-b,b}\right)^{-1}.
\label{eq:gauss-marginal-block}
\end{equation}
(ii) The difference \(E^\top\Lambda E-(E^\top\Lambda^{-1}E)^{-1}\) is positive
semidefinite of rank \(\operatorname{rank}(E^\top\Lambda F)\), and it vanishes
exactly when \(E^\top\Lambda F=0\), that is when \(\operatorname{range}E\) and
\(\operatorname{range}F\) are \(\Lambda\)-orthogonal. (iii)
\((E^\top\Lambda E)^{-1}\preceq E^\top\Lambda^{-1}E\), with equality exactly when
\(E^\top\Lambda F=0\). (iv) The determinant factorizes as
\[
\det\Lambda=\det(E^\top\Lambda E)\det\bigl(F^\top\Lambda F-F^\top\Lambda E(E^\top\Lambda E)^{-1}E^\top\Lambda F\bigr)
=\frac{\det(E^\top\Lambda E)}{\det(F^\top\Lambda^{-1}F)}.
\]
(v) A zero precision block says nothing about the corresponding covariance
block:
\begin{equation}
\Lambda=\begin{pmatrix}2&-1&0\\-1&2&-1\\0&-1&2\end{pmatrix},
\qquad
\Lambda^{-1}=\frac14\begin{pmatrix}3&2&1\\2&4&2\\1&2&3\end{pmatrix},
\label{eq:gauss-precision-not-covariance}
\end{equation}
has \(\Lambda_{13}=0\) and \((\Lambda^{-1})_{13}=\tfrac14\). \status{ESTABLISHED}

\noindent\emph{Proof.}
In the basis \([E\ F]\), \(\widetilde\Lambda=[E\ F]^\top\Lambda[E\ F]\) has
blocks \(\widetilde\Lambda_{EE}=E^\top\Lambda E\),
\(\widetilde\Lambda_{EF}=E^\top\Lambda F\), \(\widetilde\Lambda_{FF}=F^\top\Lambda F\),
and orthogonality gives \(\widetilde\Lambda^{-1}=[E\ F]^\top\Lambda^{-1}[E\ F]\).
Multiplying out the factorization
\[
\widetilde\Lambda
=\begin{pmatrix}I&\widetilde\Lambda_{EF}\widetilde\Lambda_{FF}^{-1}\\0&I\end{pmatrix}
\begin{pmatrix}\widetilde\Lambda_{EE}-\widetilde\Lambda_{EF}\widetilde\Lambda_{FF}^{-1}\widetilde\Lambda_{FE}&0\\0&\widetilde\Lambda_{FF}\end{pmatrix}
\begin{pmatrix}I&0\\\widetilde\Lambda_{FF}^{-1}\widetilde\Lambda_{FE}&I\end{pmatrix}
\]
and inverting shows that the \(EE\) block of \(\widetilde\Lambda^{-1}\) is the
inverse of the Schur complement, which is (i); the coordinate case takes for
\(E\) the columns of the identity in block \(b\). The difference in (ii) is
\(M^\top\widetilde\Lambda_{FF}^{-1}M\) with \(M=\widetilde\Lambda_{FE}\) and
\(\widetilde\Lambda_{FF}^{-1}\succ0\), hence positive semidefinite with the
rank of \(M\). For (iii), if \(0\prec S\preceq T\) then
\(T^{-1/2}ST^{-1/2}\preceq I\), so \(S^{-1}\succeq T^{-1}\), with equality only
when \(S=T\) \citep[\S7.7]{HornJohnson2013}; apply it to (ii). The
factorization with the roles of \(E\) and \(F\) exchanged gives the first
formula in (iv), and (i) with the roles exchanged identifies that Schur
complement with \((F^\top\Lambda^{-1}F)^{-1}\). Part (v) is direct
multiplication. \(\square\)

The lemma is the reason this supplement works in information form: a
restriction on the recognition family constrains conditional dependence, which
is a coordinate of \(\Lambda\), while its effect on marginal covariances runs
through every path in the remaining coordinates. Its corollaries below are the
four forms in which it is used.

\corollaryheading{The block conditional covariance is dominated by the marginal}{prop:restrict-optimum-vs-marginal}
For every block \(b\),
\begin{equation}
\Lambda_{bb}^{-1}\preceq\left(\Lambda^{-1}\right)_{bb},
\label{eq:restrict-underdispersion}
\end{equation}
with equality if and only if \(\Lambda_{b,-b}=0\). The left side is the block
covariance of the optimal factorized recognition law of
\Cref{thm:restrict-exact-block-optimum}, so that law understates every
marginal block covariance of the target except in the decoupled case.
\status{ESTABLISHED}

\noindent\emph{Proof.}
\Cref{lem:gauss-schur}(iii) with \(E\) the coordinate columns of block \(b\).
\(\square\)

\corollaryheading{Marginal precision against restricted precision}{prop:restrict-marginal-vs-restricted-precision}
Let \([B\ B_\perp]\) be orthogonal with \(B\in\R^{n\times r}\), and write
\(\widetilde\Lambda=[B\ B_\perp]^\top\Lambda[B\ B_\perp]\) with blocks indexed
by \(\parallel\) and \(\perp\). The marginal precision of the directions
\(B_\perp\) is the Schur complement
\begin{equation}
\left(B_\perp^\top\Lambda^{-1}B_\perp\right)^{-1}
=\widetilde\Lambda_{\perp\perp}
-\widetilde\Lambda_{\perp\parallel}\widetilde\Lambda_{\parallel\parallel}^{-1}\widetilde\Lambda_{\parallel\perp},
\label{eq:restrict-schur-in-basis}
\end{equation}
and
\begin{equation}
B_\perp^\top\Lambda B_\perp-\left(B_\perp^\top\Lambda^{-1}B_\perp\right)^{-1}
=B_\perp^\top\Lambda B\left(B^\top\Lambda B\right)^{-1}B^\top\Lambda B_\perp
\succeq0,
\label{eq:restrict-marginal-vs-restricted}
\end{equation}
a difference of rank at most \(r\) that vanishes exactly when
\(B^\top\Lambda B_\perp=0\). \status{ESTABLISHED}

\noindent\emph{Proof.}
\Cref{lem:gauss-schur}(i) and (ii) with \(E=B_\perp\) and \(F=B\); the rank of
\(B^\top\Lambda B_\perp\) is at most \(r\). \(\square\)

\corollaryheading{Fisher monotonicity of the mean submodel}{prop:ig-pullback-vs-pushforward}
Let \(B\in\R^{n\times k}\) have orthonormal columns and let \(B_\perp\) span
the orthogonal complement. The mean submodel
\(\{\mathcal N(B\bar m,\Lambda^{-1}):\bar m\in\R^k\}\) has Fisher metric
\(B^\top\Lambda B\), the pullback of the mean metric \(\Lambda\) along
\(\bar m\mapsto B\bar m\). The deterministic channel \(y\mapsto B^\top y\)
carries it to \(\{\mathcal N(\bar m,B^\top\Lambda^{-1}B)\}\), whose Fisher
metric is the marginal precision \((B^\top\Lambda^{-1}B)^{-1}\), and
\begin{equation}
B^\top\Lambda B-\left(B^\top\Lambda^{-1}B\right)^{-1}
=\left(B^\top\Lambda B_\perp\right)\left(B_\perp^\top\Lambda B_\perp\right)^{-1}\left(B_\perp^\top\Lambda B\right)
\succeq 0.
\label{eq:ig-restriction-marginal}
\end{equation}
The right side is the Fisher loss of the channel, and it vanishes exactly when
\(B^\top y\) is sufficient for \(\bar m\), that is when
\(B^\top\Lambda B_\perp=0\). \status{ESTABLISHED}

\noindent\emph{Proof.}
The score of the submodel is \(B^\top\Lambda(y-B\bar m)\), with covariance
\(B^\top\Lambda B\); the pushed family has covariance \(B^\top\Lambda^{-1}B\),
whose inverse is its Fisher metric. \Cref{lem:gauss-schur}(i) and (ii) with
\(E=B\), \(F=B_\perp\) give \eqref{eq:ig-restriction-marginal}. Writing
\(y=BB^\top y+B_\perp B_\perp^\top y\), the conditional covariance of the score
given \(B^\top Y\) is
\(B^\top\Lambda B_\perp\operatorname{Cov}(B_\perp^\top Y\given B^\top Y)B_\perp^\top\Lambda B\),
and \(\operatorname{Cov}(B_\perp^\top Y\given B^\top Y)=(B_\perp^\top\Lambda B_\perp)^{-1}\)
by \Cref{prop:gauss-conditional-precision}(i) in the basis \([B\ B_\perp]\); so
the right side of \eqref{eq:ig-restriction-marginal} is the expected
conditional score covariance of \Cref{thm:cg-fisher-contraction}. It vanishes
exactly when \(B^\top\Lambda B_\perp=0\), which is when the score, and hence the
likelihood ratio, depends on \(y\) only through \(B^\top y\). \(\square\)

\corollaryheading{Full-rank inclusion and averaging}{cor:ig-full-rank-restriction-pushforward}
Let \(S\in\R^{n\times k}\) have full column rank, set \(M_S=S^\top S\), and
let \(S^\dagger=M_S^{-1}S^\top\). The mean submodel \(\mu=S\bar\mu\) has Fisher
matrix \(S^\top\Lambda S\), while the mean of the pushforward under
\(y\mapsto S^\dagger y\) has Fisher matrix
\begin{equation}
\Lambda_{\rm av}
=M_S\left(S^\top\Lambda^{-1}S\right)^{-1}M_S.
\label{eq:ig-full-rank-average-precision}
\end{equation}
With \(B=SM_S^{-1/2}\) completed to an orthogonal matrix \([B\ B_\perp]\),
\begin{equation}
S^\top\Lambda S-\Lambda_{\rm av}
=M_S^{1/2}\left[
B^\top\Lambda B-\left(B^\top\Lambda^{-1}B\right)^{-1}
\right]M_S^{1/2}\succeq0,
\label{eq:ig-full-rank-restriction-marginal}
\end{equation}
with equality exactly when \(\operatorname{range}S\) and its Euclidean
orthogonal complement are \(\Lambda\)-orthogonal. \status{ESTABLISHED}

\noindent\emph{Proof.}
The pushforward covariance is
\(S^\dagger\Lambda^{-1}(S^\dagger)^\top=M_S^{-1}S^\top\Lambda^{-1}SM_S^{-1}\),
whose inverse is \eqref{eq:ig-full-rank-average-precision}. Substituting
\(S=BM_S^{1/2}\) into both matrices gives
\eqref{eq:ig-full-rank-restriction-marginal}, and
\Cref{prop:ig-pullback-vs-pushforward} supplies its positivity and equality
condition. \(\square\)

\section{Moments and vector residuals}

\propositionheading{Finite-second-moment congruence}{prop:geo-moment-pushforward}
Let a belief law \(q\) on \(\R^K\) have finite mean \(\mu\) and covariance
\(\Sigma\), and let \(g\in G\) act by the linear map \(\rho_b(g)\). The pushed
law \(\rho_b(g)_\#q\) has
\begin{equation}
\mu'=\rho_b(g)\mu,
\qquad
\Sigma'=\rho_b(g)\Sigma\rho_b(g)^\top,
\label{eq:geo-moment-pushforward}
\end{equation}
and if \(\Sigma\succ0\) its precision is
\begin{equation}
\Lambda'=\rho_b(g)^{-\top}\Lambda\rho_b(g)^{-1},
\qquad\Lambda=\Sigma^{-1}.
\label{eq:geo-precision-inverse-congruence}
\end{equation}
The model channel has the same statement in dimension \(d_m\).
\status{ESTABLISHED}

\noindent\emph{Proof.}
Linearity of expectation gives the mean, and the centered pushed variable
\(\rho_b(g)(Y-\mu)\) has covariance \(\rho_b(g)\Sigma\rho_b(g)^\top\);
inversion gives \eqref{eq:geo-precision-inverse-congruence}. \(\square\)

\subsection{Residuals of aligned and twisted links}
\label{sec:geo-flatness}

Let \(\mu_i\) be a belief-channel vector coordinate of agent \(i\), let
\(U_i\) be the frame of \(u_i\) relative to a reference point of the fiber, as
in \Cref{prop:geo-edge-twist-invariance}, and let \(e:j\to i\) be an edge copy
with belief link \(\Theta^b_e\) and per-edge twist
\(\mathring\Theta^b_e=U_i^{-1}\Theta^b_eU_j\) (\eqref{eq:geo-edge-twist}).

\propositionheading{Vector residuals}{prop:gauss-vector-residuals}
With \(z_i=\rho_b(U_i)^{-1}\mu_i\),
\begin{equation}
\mu_i-\rho_b(\Theta_e^b)\mu_j
=\rho_b(U_i)\bigl(z_i-\rho_b(\mathring\Theta_e^b)z_j\bigr),
\label{eq:geo-regime-two-residual}
\end{equation}
and for frame-aligned links (\Cref{hyp:geo-flat-links}), whose twists equal
\(e_G\),
\begin{equation}
\mu_i-\rho_b(\Theta_e^b)\mu_j
=\rho_b(U_i)(z_i-z_j).
\label{eq:geo-trivialized-residual}
\end{equation}
More generally, if \(\Theta_e^b=U'_i(U'_j)^{-1}\) on a belief-trivializing
subfamily of links and \(z'_i=\rho_b(U'_i)^{-1}\mu_i\), then
\begin{equation}
\mu_i-\rho_b(\Theta_e^b)\mu_j
=\rho_b(U'_i)(z'_i-z'_j).
\label{eq:geo-subfamily-trivialized-residual}
\end{equation}
\status{ESTABLISHED}

\noindent\emph{Proof.}
Substitute \(\mu_i=\rho_b(U_i)z_i\) and
\(\rho_b(\Theta^b_e)=\rho_b(U_i)\rho_b(\mathring\Theta^b_e)\rho_b(U_j)^{-1}\);
\eqref{eq:geo-trivialized-residual} is the case \(\mathring\Theta^b_e=e_G\), and
\eqref{eq:geo-subfamily-trivialized-residual} is the same substitution with the
potentials \(U'_i\). \(\square\)

Belief-link trivialization says nothing about model links or either cross
morphism.

\section{The directed linear-Gaussian model}
\label{sec:gen-gaussian}

\definitionheading{Directed linear-Gaussian forest}{def:gen-lg-forest}
Take
\begin{equation}
k_i\in\R^K,
\qquad
m_i\in\R^{d_m},
\qquad
o_i\in\R^{d_{o,i}},
\label{eq:gen-lg-domains}
\end{equation}
and a directed forest on the agents, so that each nonroot \(i\) has one parent
\(j=\operatorname{par}(i)\), joined by an edge copy \(j\to i\) with links
\(\Theta^b_{ij},\Theta^m_{ij}\in G\) (\Cref{def:geo-graph-links}). For a
nonroot,
\begin{align}
m_i\given m_j,X
&\sim\mathcal N\left(
F_i^m\rho_m(\Theta_{ij}^m)m_j+\xi_i^m,
\Sigma_i^m\right),
\label{eq:gen-lg-model}\\
k_i\given k_j,m_i,X
&\sim\mathcal N\left(
F_i^b\rho_b(\Theta_{ij}^b)k_j+\mathsf B_im_i+\xi_i^b,
\Sigma_i^b\right),
\label{eq:gen-lg-state}\\
o_i\given k_i,m_i,X
&\sim\mathcal N\left(
\mathsf C_i^bk_i+\mathsf C_i^mm_i+\xi_i^o,
\Sigma_i^o\right),
\label{eq:gen-lg-obs}
\end{align}
with parameter shapes
\begin{align}
&\rho_b(\Theta_{ij}^b),F_i^b\in\R^{K\times K},
&&\rho_m(\Theta_{ij}^m),F_i^m\in\R^{d_m\times d_m},
&&\mathsf B_i\in\R^{K\times d_m},
\label{eq:gen-lg-shapes-latent}\\
&\mathsf C_i^b\in\R^{d_{o,i}\times K},
&&\mathsf C_i^m\in\R^{d_{o,i}\times d_m},
&&\xi_i^b\in\R^K,
\quad
\xi_i^m\in\R^{d_m},
\quad
\xi_i^o\in\R^{d_{o,i}}.
\label{eq:gen-lg-shapes-obs}
\end{align}
Each root \(r\) has the laws
\begin{align}
m_r\given X
&\sim\mathcal N(\bar m_r,\Sigma_{r,0}^m),
\label{eq:gen-lg-root-model}\\
k_r\given m_r,X
&\sim\mathcal N(\mathsf B_rm_r+\xi_{r,0}^b,\Sigma_{r,0}^b),
\label{eq:gen-lg-root-state}
\end{align}
followed by \eqref{eq:gen-lg-obs}. A general directed acyclic graph replaces
each parent term by a finite sum. The represented links are invertible; the
gains \(F_i^b,F_i^m\) need not be.

\hypothesisheading{Positive-definite receiving covariances}{hyp:gen-lg-spd}
Every receiving covariance is positive definite:
\begin{equation}
\Sigma_i^m,\Sigma_{r,0}^m\in\Sym_{++}^{d_m},
\qquad
\Sigma_i^b,\Sigma_{r,0}^b\in\Sym_{++}^{K},
\qquad
\Sigma_i^o\in\Sym_{++}^{d_{o,i}}.
\label{eq:gen-lg-spd}
\end{equation}
The hypothesis excludes degenerate Gaussian kernels, whose laws have no
Lebesgue density on the receiving space.

\propositionheading{Normalization and root moments}{prop:gen-lg-normalization}
Under \Cref{hyp:gen-lg-spd}, every kernel of \Cref{def:gen-lg-forest} is a
nondegenerate normalized density and the assembled directed law is a
probability measure. The root belief coordinate has
\begin{equation}
\bar k_r=\mathsf B_r\bar m_r+\xi_{r,0}^b,
\qquad
\Sigma_{r,0}^{b,\mathrm{marg}}
=\Sigma_{r,0}^b+\mathsf B_r\Sigma_{r,0}^m\mathsf B_r^\top\succeq\Sigma^b_{r,0},
\label{eq:gen-lg-root-marginal}
\end{equation}
and, for a root model law \(s=\mathcal N(\mu_s,\Sigma_s)\), the law map induced
by the root kernel is
\begin{equation}
\widetilde\phi_r(s)
=\mathcal N\left(
\mathsf B_r\mu_s+\xi_{r,0}^b,\
\Sigma_{r,0}^b+\mathsf B_r\Sigma_s\mathsf B_r^\top
\right).
\label{eq:gen-lg-induced-law-map}
\end{equation}
This law map, of which \(\mathsf B_r\) is one parameter, is the local realization of
\(\widetilde\Phi\). \status{ESTABLISHED}

\noindent\emph{Proof.}
A Gaussian density with positive-definite covariance integrates to one for
every mean, so \Cref{prop:gen-exact-normalization} applies. The tower property
gives the means and the law of total covariance gives the covariances; the
added term \(\mathsf B_r\Sigma_s\mathsf B_r^\top\) is positive semidefinite. \(\square\)

At fixed model latents, write \(F_i:=F_i^b\), \(\Theta_{ij}:=\rho_b(\Theta^b_{ij})\in\Gamma_b\),
and \(\Lambda_i:=(\Sigma_i^b)^{-1}\); the belief sector is the forest of
receiving kernels
\begin{equation}
k_i\given k_j\sim\mathcal N\left(F_i\Theta_{ij}k_j,\Lambda_i^{-1}\right),
\qquad
j=\operatorname{par}(i),
\label{eq:gauss-receiving-kernel}
\end{equation}
with a proper prior \(\mathcal N(0,\Lambda_r^{-1})\) at each root; the bridge
term, the offsets, and a nonzero prior mean shift only the information vector.
A single subscript on \(\Lambda\) denotes the innovation precision of one
receiving kernel, a double subscript a block of the assembled precision.

\propositionheading{Assembled precision blocks}{prop:gauss-assembled-precision-blocks}
Write \(\Omega_i:=F_i\Theta_{i\operatorname{par}(i)}\) for the transition of a
nonroot \(i\) and \(\operatorname{ch}(i)=\{i':\operatorname{par}(i')=i\}\). The
joint precision of the belief sector has blocks
\begin{equation}
\Lambda_{i\operatorname{par}(i)}=-\Lambda_i\Omega_i,
\qquad
\Lambda_{ii}=\Lambda_i+\sum_{i'\in\operatorname{ch}(i)}\Omega_{i'}^\top\Lambda_{i'}\Omega_{i'},
\label{eq:gauss-directed-blocks}
\end{equation}
and \(\Lambda_{ij}=0\) whenever \(i\) and \(j\) are not adjacent in the forest.
In the joint precision of \((k,m)\), the block between \(k_i\) and \(m_i\) is
\(-\Lambda_i\mathsf B_i+\mathsf C_i^{b\top}(\Sigma_i^o)^{-1}\mathsf C_i^m\in\R^{K\times d_m}\), whose
first term is the prior contribution; a family built from square
\(K\times K\) agent blocks cannot host it when \(K\neq d_m\).
\status{ESTABLISHED}

\noindent\emph{Proof.}
The negative log joint density of the belief sector is
\(\tfrac12\sum_i(k_i-\Omega_ik_{\operatorname{par}(i)})^\top\Lambda_i(k_i-\Omega_ik_{\operatorname{par}(i)})\)
plus terms free of \(k\), with no parent term at a root. One summand
contributes \(\Lambda_i\) to the \(ii\) block, \(-\Lambda_i\Omega_i\) to the
\(i\operatorname{par}(i)\) block, and \(\Omega_i^\top\Lambda_i\Omega_i\) to the
parent's diagonal block; collecting them gives \eqref{eq:gauss-directed-blocks}.
In the joint energy, the factors \eqref{eq:gen-lg-state} and
\eqref{eq:gen-lg-obs} are the only ones containing both \(k_i\) and \(m_i\),
and their mixed second derivatives are \(-\Lambda_i\mathsf B_i\) and
\(\mathsf C_i^{b\top}(\Sigma_i^o)^{-1}\mathsf C_i^m\). \(\square\)

\propositionheading{Gaussian kernel covariance}{prop:gen-lg-kernel-covariance}
Let \(g_i\in G\) act by \eqref{eq:gen-gaussian-diagonal-action}, write
\(\rho^b_i:=\rho_b(g_i)\), \(\rho^m_i:=\rho_m(g_i)\), let the links transform by
\eqref{eq:gen-gauge-links}, so that
\(\rho_x(\Theta^{x\prime}_{ij})=\rho^x_i\rho_x(\Theta^x_{ij})(\rho^x_j)^{-1}\),
and let the parameters transform by
\begin{align}
F_i^{b\prime}
&=\rho^b_iF_i^b(\rho^b_i)^{-1},
&
F_i^{m\prime}
&=\rho^m_iF_i^m(\rho^m_i)^{-1},
\label{eq:gen-gauge-selfmaps}\\
\mathsf B_i'
&=\rho^b_i\mathsf B_i(\rho^m_i)^{-1},
&
\xi_i^{b\prime}
&=\rho^b_i\xi_i^b,
\qquad
\xi_i^{m\prime}=\rho^m_i\xi_i^m,
\label{eq:gen-gauge-bridge-offsets}\\
\Sigma_i^{b\prime}&=\rho^b_i\Sigma_i^b(\rho^b_i)^\top,
&
\Sigma_i^{m\prime}&=\rho^m_i\Sigma_i^m(\rho^m_i)^\top,
\label{eq:gen-gauge-covariances}
\end{align}
so that for either channel
\begin{equation}
\Lambda_i=\Sigma_i^{-1},
\qquad
\Sigma_i'=\rho_i\Sigma_i\rho_i^\top,
\qquad
\Lambda_i'=\rho_i^{-\top}\Lambda_i\rho_i^{-1},
\label{eq:gen-gauge-precision}
\end{equation}
with the root laws
\begin{align}
\bar m_r'
&=\rho^m_r\bar m_r,
&
\Sigma_{r,0}^{m\prime}
&=\rho^m_r\Sigma_{r,0}^m(\rho^m_r)^\top,
\label{eq:gen-gauge-root-model}\\
\mathsf B_r'
&=\rho^b_r\mathsf B_r(\rho^m_r)^{-1},
&
\xi_{r,0}^{b\prime}
&=\rho^b_r\xi_{r,0}^b,
\qquad
\Sigma_{r,0}^{b\prime}
=\rho^b_r\Sigma_{r,0}^b(\rho^b_r)^\top,
\label{eq:gen-gauge-root-state}
\end{align}
and gauge-inert observations
\begin{equation}
\mathsf C_i^{b\prime}=\mathsf C_i^b(\rho^b_i)^{-1},
\qquad
\mathsf C_i^{m\prime}=\mathsf C_i^m(\rho^m_i)^{-1},
\qquad
\xi_i^{o\prime}=\xi_i^o,
\qquad
\Sigma_i^{o\prime}=\Sigma_i^o.
\label{eq:gen-gauge-observation}
\end{equation}
Then every transformed kernel is the receiving-frame pushforward of the
original, so the kernel identities \eqref{eq:gen-gauge-pushforward-model},
\eqref{eq:gen-gauge-pushforward-state}, and \eqref{eq:gen-gauge-record-kernel}
hold and the joint law and its evidence obey
\Cref{prop:gen-product-evidence-invariance}. The induced root law map
\eqref{eq:gen-lg-induced-law-map} obeys the frame law
\eqref{eq:geo-tildephi-gauge-law} of \(\widetilde\Phi\), the root marginal
\eqref{eq:gen-lg-root-marginal} transforms by belief-channel congruence, and
positive definiteness is preserved in each channel. \status{ESTABLISHED}

\noindent\emph{Proof.}
For the belief kernel, the transformed conditional mean at transformed inputs is
\begin{equation}
F_i^{b\prime}\rho_b(\Theta_{ij}^{b\prime})(\rho^b_jk_j)
+\mathsf B_i'(\rho^m_im_i)+\xi_i^{b\prime}
=\rho^b_i\left(
F_i^b\rho_b(\Theta_{ij}^b)k_j+\mathsf B_im_i+\xi_i^b
\right),
\label{eq:gen-gauge-mean-identity}
\end{equation}
and the covariance transforms by \eqref{eq:gen-gauge-covariances}, so the
transformed Gaussian is the pushforward of the original under \(\rho^b_i\). The
model kernel is the same computation without a bridge. At the observation
kernel,
\begin{equation}
\mathsf C_i^{b\prime}(\rho^b_ik_i)+\mathsf C_i^{m\prime}(\rho^m_im_i)+\xi_i^{o\prime}
=\mathsf C_i^bk_i+\mathsf C_i^mm_i+\xi_i^o,
\label{eq:gen-gauge-observation-identity}
\end{equation}
with fixed covariance. The root calculation is identical, and the law map
\eqref{eq:gen-lg-induced-law-map} of a pushed model law
\(\rho^m_{r\#}s\) is \(\rho^b_{r\#}\widetilde\phi_r(s)\) by the same
substitution. Composition is \Cref{prop:gen-product-evidence-invariance}.
Both summands of \eqref{eq:gen-lg-root-marginal} transform by belief-channel
congruence, and congruence by an invertible matrix preserves positive
definiteness. \(\square\)

\section{The Gaussian M-coordinate}

\propositionheading{Gaussian M-coordinate and its existence proviso}{prop:gen-lg-m-step}
Suppose the factors whose parameters lie in \(\theta\) carry disjoint
parameter blocks, that each such factor is a law \(\mathcal N(Fw,\Sigma)\) of
its target coordinate \(t\) given a regressor vector \(w\) listing its parent
coordinates and the constant \(1\), with \(F\) free over all matrices of its
shape and \(\Sigma\) in the open cone of \eqref{eq:gen-lg-spd}, and that
\(\mathbb Q_{V,o,X}\) has finite second moments. The expected complete
log-likelihood of \eqref{eq:elbo-m-coordinate} is then a sum of separately
maximized factor terms. For one factor, with expectations under
\(\mathbb Q_{V,o,X}\), put \(M_{ww}=\E[ww^\top]\), let \(F_*\) solve
\(F_*M_{ww}=\E[tw^\top]\), and set \(\Sigma_*=\E[(t-F_*w)(t-F_*w)^\top]\). Then
\begin{equation}
\E\bigl[(t-Fw)(t-Fw)^\top\bigr]=\Sigma_*+(F-F_*)M_{ww}(F-F_*)^\top,
\qquad
\E\log\mathcal N(t;Fw,\Sigma)\leq-\tfrac12\log\det(2\pi e\Sigma_*),
\label{eq:elbo-gaussian-m-step}
\end{equation}
with equality at \((F_*,\Sigma_*)\) when \(\Sigma_*\succ0\). If \(\Sigma_*\) is
singular the term is unbounded above and has no maximizer. The matrix
\(\Sigma_*\) is positive definite exactly when no nonzero linear combination of
\(t\) is almost surely an affine function of the parents, which holds for a
latent target whenever \(\mathbb Q_{V,o,X}\) has a Lebesgue density on the
latent coordinates. Two cases violate the proviso: a record factor at its
single observed value, whose intercept can equal that value with zero
coefficients on the parents, giving \(\Sigma_*=0\); and a link learned inside
the invertible group with the rest of its mean map fixed, which for \(K=1\),
link map \(\mathsf t=\rho_b(\Theta^b_{ij})\neq0\), and \(k_i,k_j\) independent
and centered, has covariance-optimized term
\(-\tfrac12\log\bigl(2\pi e(\E k_i^2+\mathsf t^2\E k_j^2)\bigr)\), which
approaches its supremum \(-\tfrac12\log(2\pi e\E k_i^2)\) as
\(\mathsf t\to0\) without attaining it. \status{ESTABLISHED}

\noindent\emph{Proof.}
Expanding \(t-Fw=(t-F_*w)-(F-F_*)w\) and using the normal equations to cancel
the cross term gives the first identity. For \(\Sigma_*\succ0\) and
\(\varrho_k\) the eigenvalues of \(\Sigma^{-1/2}\Sigma_*\Sigma^{-1/2}\),
\(-\log\det\Sigma-\Tr(\Sigma^{-1}\Sigma_*)
=-\log\det\Sigma_*+\sum_k(\log\varrho_k-\varrho_k)\), which is maximal exactly
at \(\Sigma=\Sigma_*\); the first identity orders every residual second moment
above \(\Sigma_*\), and \(\Tr(\Sigma^{-1}\cdot)\) is monotone, which gives the
inequality. If \(\Sigma_*\) is singular, \(\Sigma=\Sigma_*+\varepsilon I\) with
\(\varepsilon\downarrow0\) drives the term to \(+\infty\). A nonzero \(a\) with
\(a^\top\Sigma_*a=0\) makes \(a^\top t\) almost surely affine in the parents,
and such a relation confines a latent vector to a Lebesgue-null hyperplane.
The two witnesses are direct evaluation. \(\square\)

\section{Gaussian selective communication}

The comparison maps \(S^Q_e,S^K_e\), the value map \(\mathsf L^{\mathrm{val}}_e\),
and the encoder, transmission, and decoder maps
\(\mathsf L^{\mathrm{enc}}_e,\mathsf L^{\mathrm{tr}}_e,\mathsf L^{\mathrm{dec}}_e\)
are those of \Cref{sec:selective-comparison,sec:selective-records}.

\propositionheading{A comparison lift requires receiver geometry}{prop:geo-selective-lift}
Fix \(q_i=\mathcal N(\mu_i,\Sigma_i)\) with \(\Sigma_i\succ0\) and a full-row-rank
\(S^Q_e\). (i) The reduced covariance \(S^Q_e\Sigma_iS_e^{Q\top}\) is positive
definite. (ii) The lift
\begin{equation}
R_e=\Sigma_iS_e^{Q\top}(S^Q_e\Sigma_iS_e^{Q\top})^{-1},\qquad
S^Q_eR_e=I,\qquad \mathsf L^{\mathrm{cmp}}_e=R_eS^K_e,
\label{eq:geo-selective-gaussian-lift}
\end{equation}
makes \(R_ey\) the unique minimizer of \(\delta^\top\Sigma_i^{-1}\delta\) subject
to \(S^Q_e\delta=y\), and
\(\E[k_i\given S^Q_ek_i=y]=\mu_i+R_e(y-S^Q_e\mu_i)\). (iii) If comparison and
content share a Gaussian source \(z\sim\mathcal N(\mu_z,\Sigma_z)\),
\(\Sigma_z\succ0\), with content map \(\mathsf L\) and full-row-rank key map
\(S^K_e\), then
\begin{equation}
\E[\mathsf Lz\given S^K_ez=y]=\mathsf L\mu_z+
\mathsf L\Sigma_zS_e^{K\top}(S^K_e\Sigma_zS_e^{K\top})^{-1}(y-S^K_e\mu_z),
\label{eq:obs-selective-shared-latent-content}
\end{equation}
and the correction vanishes for all \(y\) exactly when
\(\mathsf L\Sigma_zS_e^{K\top}=0\). (iv) Under \(S_e^{Q\prime}=\phi^{\mathrm{cmp}}_eS^Q_e\rho_b(g_i)^{-1}\)
with \(\phi^{\mathrm{cmp}}_e\) invertible on the comparison space, and
\(\Sigma_i'=\rho_b(g_i)\Sigma_i\rho_b(g_i)^\top\), the lift transforms by
\(R_e'=\rho_b(g_i)R_e(\phi^{\mathrm{cmp}}_e)^{-1}\). (v) In the comparison-record model of
\Cref{prop:obs-selective-receiver-update}, with \(S^Q_i\) of full row rank,
\(R_e\) formed from \(S^Q_i\), and \(\widetilde W_i^{-1}\to0\),
\(\mathsf L^{\mathrm{val}}_e\to R_eS^K_e\) and
\(\Sigma_i^+\to\Sigma_i-R_eS^Q_i\Sigma_i\). (vi) For comparison and content
copies \(Y=V=z\sim\mathcal N(0,1)\), every product law \(q(Y)q(V)\) with
continuous marginals has infinite relative entropy to the diagonal-supported
joint, even with identical comparison marginals. \status{ESTABLISHED}

\noindent\emph{Proof.}
(i) For \(u\neq0\), \(S_e^{Q\top}u\neq0\) by full row rank, so
\(u^\top S^Q_e\Sigma_iS_e^{Q\top}u>0\). (ii) Write a feasible displacement as
\(R_ey+w\) with \(S^Q_ew=0\); its squared norm has cross term
\(y^\top(S^Q_e\Sigma_iS_e^{Q\top})^{-1}S^Q_ew=0\) and the additional term
\(w^\top\Sigma_i^{-1}w\), positive unless \(w=0\). The residual
\(k_i-\mu_i-R_e(S^Q_ek_i-S^Q_e\mu_i)\) has zero covariance with \(S^Q_ek_i\),
hence is independent of it, which gives the conditional mean. (iii) The same
argument applied to \(z\) and followed by \(\mathsf L\). (iv) Substitution.
(v) The Woodbury identity
\(\Sigma_i^+S_i^{Q\top}\widetilde W_i=\Sigma_iS_i^{Q\top}(\Sigma_i^{\rm pred})^{-1}\),
with \(\Sigma_i^{\rm pred}=\widetilde W_i^{-1}+S^Q_i\Sigma_iS_i^{Q\top}\) of
\eqref{eq:obs-selective-receiver-covariances}, gives
\(\mathsf L^{\mathrm{val}}_e=\Sigma_iS_i^{Q\top}(\Sigma_i^{\rm pred})^{-1}S^K_e\)
and \(\Sigma_i^+=\Sigma_i-\Sigma_iS_i^{Q\top}(\Sigma_i^{\rm pred})^{-1}S^Q_i\Sigma_i\);
as \(\widetilde W_i^{-1}\to0\), \((\Sigma_i^{\rm pred})^{-1}\to(S^Q_i\Sigma_iS_i^{Q\top})^{-1}\)
by (i). (vi) The joint of \((Y,V)\) is concentrated on the diagonal, which a
product of continuous laws assigns mass zero, so absolute continuity fails.
\(\square\)

The content map \(\mathsf L_e\) of a communication layer is a separate
parameter; the lift \(\mathsf L^{\mathrm{cmp}}_e\) is what receiver geometry
supplies when content is read back from a comparison.

\propositionheading{The linear Gaussian record}{prop:obs-selective-gaussian-record}
Let the three stages of \eqref{eq:obs-selective-kernel-composition} be
\(h^-=\mathsf L^{\mathrm{enc}}_ek_j+\epsilon_-\),
\(h^+=\mathsf L^{\mathrm{tr}}_eh^-+\epsilon_+\), and
\(o_e=\mathsf L^{\mathrm{dec}}_eh^++\epsilon_o\), with mutually independent
centered Gaussian noises, independent of \(k_j\), of covariances
\(\Sigma^-_e,\Sigma^+_e,\Sigma^o_e\succeq0\). (i) Then
\begin{equation}
o_e\given k_j\sim\mathcal N(\mathsf L_ek_j,\Sigma^{\mathrm{rec}}_e),
\qquad
\Sigma^{\mathrm{rec}}_e=\mathsf L^{\mathrm{dec}}_e\mathsf L^{\mathrm{tr}}_e\Sigma^-_e\mathsf L_e^{\mathrm{tr}\top}\mathsf L_e^{\mathrm{dec}\top}
+\mathsf L^{\mathrm{dec}}_e\Sigma^+_e\mathsf L_e^{\mathrm{dec}\top}+\Sigma^o_e,
\label{eq:obs-selective-gaussian-record}
\end{equation}
with \(\mathsf L_e=\mathsf L^{\mathrm{dec}}_e\mathsf L^{\mathrm{tr}}_e\mathsf L^{\mathrm{enc}}_e\);
the kernel is normalized at every rank, has a Lebesgue density when
\(\Sigma^{\mathrm{rec}}_e\succ0\), for instance when \(\Sigma^o_e\succ0\), and
for input \(\mathcal N(\mu_j,\Sigma_j)\) gives
\(o_e\sim\mathcal N(\mathsf L_e\mu_j,\Sigma^{\mathrm{rec}}_e+\mathsf L_e\Sigma_j\mathsf L_e^\top)\);
if the receiver frame changes by \(g_i\), with
\(\mathsf L^{\mathrm{dec}\prime}_e=\rho_b(g_i)\mathsf L^{\mathrm{dec}}_e\) and
\(\Sigma^{o\prime}_e=\rho_b(g_i)\Sigma^o_e\rho_b(g_i)^\top\), then
\(\Sigma^{\mathrm{rec}\prime}_e=\rho_b(g_i)\Sigma^{\mathrm{rec}}_e\rho_b(g_i)^\top\).
Deterministic placement is normalized as a measure but need not have a
Lebesgue density. (ii) For \(\Sigma^{\mathrm{rec}}_e\succ0\) and
\(k_j\sim\mathcal N(0,\Sigma_j)\),
\begin{equation}
I(k_j;o_e)=\tfrac12\log
\frac{\det(\Sigma^{\mathrm{rec}}_e+\mathsf L_e\Sigma_j\mathsf L_e^\top)}{\det\Sigma^{\mathrm{rec}}_e},
\label{eq:obs-selective-gaussian-information}
\end{equation}
which is unchanged by a congruent transformation of signal and noise; for
\(\mathsf L_e=(1,0)\), \(\Sigma_j=I_2\), and scalar noise variance
\(\sigma^2>0\) it equals \(\tfrac12\log(1+1/\sigma^2)\), so at the fixed rank
one it takes every value in \((0,\infty)\). (iii) The comparison energy has
the normalized realization
\(o_e\given k_i,k_j\sim\mathcal N(S^Q_ek_i-S^K_ek_j,W_e^{-1})\), \(W_e\succ0\),
whose negative log density at \(o_e=0\) is \(V_e\) of
\eqref{eq:obs-selective-response} plus the Gaussian normalizer; with a proper
normalized baseline and finitely many such records, the evidence at this slice
is finite and positive, for arbitrary rectangular \(S^Q_e,S^K_e\).
\status{ESTABLISHED}

\noindent\emph{Proof.}
(i) Substituting the stages gives
\(o_e=\mathsf L_ek_j+\mathsf L^{\mathrm{dec}}_e\mathsf L^{\mathrm{tr}}_e\epsilon_-+\mathsf L^{\mathrm{dec}}_e\epsilon_++\epsilon_o\),
a sum of independent Gaussians, so the conditional covariance is
\(\Sigma^{\mathrm{rec}}_e\) and the marginal adds \(\mathsf L_e\Sigma_j\mathsf L_e^\top\).
A Gaussian law is a probability measure for every positive-semidefinite
covariance. (ii) The difference of the output entropy and the conditional
entropy \(\tfrac12\log\det(2\pi e\Sigma^{\mathrm{rec}}_e)\) is the displayed
ratio, and a common congruence multiplies both determinants by the same
factor; the example is direct evaluation. (iii) Each likelihood is bounded by
its normalizer and strictly positive, so integrating it against a normalized
baseline gives finite positive evidence. \(\square\)
